
\section{拉格朗日弗洛尔上同调的一般理论}
\label{sec:1-2}

从现在起，我们对弗洛尔理论采用上同调记号。因此，除非另有说明，我们只使用上同调而不用同调。

为了完整刻画拉格朗日子流形的弗洛尔上同调，必须使用滤过 $A_\infty$ 代数和滤过
$A_\infty$ 模的语言。我们将在第 3--5 章中系统地叙述这些概念；这里只给出陈述下文定理 A
所需的定义。

首先引入系数环 $\Lambda_{0,\mathrm{nov}}$：
\[
  \Lambda_{0,\mathrm{nov}}
  =\left\{
    \left.
    \sum_{i=0}^{\infty}a_iT^{\lambda_i}e^{n_i}
    \ \right|\ 
    a_i\in\mathbb Q,
    \ \lambda_i\in\mathbb R_{\geq 0},
    \ n_i\in\mathbb Z,
    \ \lim_{i\to\infty}\lambda_i=\infty
  \right\}.
\]
这里 $T$ 和 $e$ 是形式生成元，$\deg e=2$，$\deg T=0$。我们称
$\Lambda_{0,\mathrm{nov}}$ 为万有诺维科夫环\liftIndex{I024}。可以假设 $\lambda_0=0$。

考虑一个（$\mathbb Q$ 上的）分次向量空间 $C$，以及由
$C[1]^k=C^{k+1}$ 定义的移位空间\liftIndex{I028} $C[1]$。考虑一系列运算
\[
  \mathfrak m_{k,i}\colon
  \underbrace{C[1]\otimes\cdots\otimes C[1]}_{k}
  \longrightarrow C[1],
\]
其次数为 $1-\mu_i$。这里 $\mu_i$ 是偶数，且 $\mu_0=0$。对每个
$k=0,1,2,\ldots$，令
\[
  \mathfrak m_k=\sum_{i=0}^{\infty}T^{\lambda_i}e^{\mu_i/2}\mathfrak m_{k,i}.
\]
这里 $\lambda_i$ 是非负实数，满足 $\lambda_0=0$，当 $i>0$ 时 $\lambda_i>0$，并且
$\lim_{i\to\infty}\lambda_i=\infty$。我们假设 $\mathfrak m_{0,0}=0$。这些运算
$\mathfrak m_k$ 定义同态
\[
  \mathfrak m_k\colon
  \underbrace{(C[1]\widehat\otimes\Lambda_{0,\mathrm{nov}})
  \otimes\cdots\otimes
  (C[1]\widehat\otimes\Lambda_{0,\mathrm{nov}})}_{k}
  \longrightarrow C[1]\widehat\otimes\Lambda_{0,\mathrm{nov}}.
\]
这里以及下文中，$\widehat\otimes$ 表示代数张量积关于适当滤过的完备化（见第 3.1、3.2 节）。

如果这些运算 $\mathfrak m_k$ 满足二次关系
\begin{equation}
  \sum_{k_1+k_2=k+1}\sum_{j=1}^{k_1}
  (-1)^{*}\mathfrak m_{k_1}
  \bigl(x_1,\ldots,
  \mathfrak m_{k_2}(x_j,\ldots,x_{j+k_2-1}),
  \ldots,x_k\bigr)=0,
  \tag{1.4}\label{eq:1-4}
\end{equation}
其中 $k=0,1,2,\ldots$，且
$*=\deg x_1+\cdots+\deg x_{j-1}+j-1$，我们就说它们在
$C[1]\widehat\otimes\Lambda_{0,\mathrm{nov}}$ 上定义了一个滤过 $A_\infty$ 代数结构\liftIndex{I002}。
我们称~\eqref{eq:1-4} 为 $A_\infty$ 关系或 $A_\infty$ 公式\liftIndex{I003}（见定义 3.2.20 和注 3.2.21）。

\begin{remark}{1.5}\label{rem:1-5}
环 $\Lambda_{0,\mathrm{nov}}$ 是谢尔盖·诺维科夫（Sergei P. Novikov）在其闭
1-形式的莫尔斯理论中引入的环的一种版本 [Nov81]，而该理论与弗洛尔上同调密切相关。因此，
本书称它为（万有）诺维科夫环。（在诺维科夫之前，类似的环已经在若干不同背景中出现过。）
弗洛尔最先注意到，诺维科夫环是弗洛尔上同调应当使用的正确
系数环 [Flo89I]。第四位署名作者小野薫（Kaoru Ono）[Ono95] 与霍费尔、
萨拉蒙 [HoSa95] 分别独立地在半正辛流形上
证明阿诺尔德猜想~(1.1) 时使用了诺维科夫环。
\end{remark}

为了在有理数系数上研究一般的弗洛尔理论，我们需要假设拉格朗日子流形
$L_i$（$i=0,1$）是相对 spin 的\liftIndex{I026}。

\begin{definition}{1.6}\label{def:1-6}
若子流形 $L\subset M$ 可定向，并且存在一个类
$st\in H^2(M;\mathbb Z_2)$，其在 $L$ 上的限制等于 $L$ 的第二施蒂费尔--惠特尼类
$w_2(L)$，则称 $L$ 是相对 spin 的。

若 $M$ 的一对子流形 $(L_1,L_0)$ 都可定向，并且存在一个类
$st\in H^2(M;\mathbb Z_2)$，其在两个 $L_i$（$i=0,1$）上的限制分别等于
$w_2(L_i)$，则称 $(L_1,L_0)$ 是相对 spin 对\liftIndex{I027}。
\end{definition}

令
\[
  H^*(L;\Lambda_{0,\mathrm{nov}})
  =H^*(L;\mathbb Q)\otimes_{\mathbb Q}\Lambda_{0,\mathrm{nov}}.
\]

\begin{theorem}{A}\label{thm:A}
对每个相对 spin 拉格朗日子流形 $L$，都可以在
$H^*(L;\Lambda_{0,\mathrm{nov}})$ 上赋予一个滤过 $A_\infty$ 代数结构
$\mathfrak m_k$；该结构在同构意义下良定义。

若 $\psi\colon(M,L)\to(M',L')$ 是辛微分同胚，则可为它赋予一个滤过
$A_\infty$ 代数同构
\[
  \psi_*\coloneqq(\psi^{-1})^*\colon
  H^*(L;\Lambda_{0,\mathrm{nov}})
  \longrightarrow H^*(L';\Lambda_{0,\mathrm{nov}}),
\]
其同伦类仅依赖于辛微分同胚
$\psi\colon(M,L)\to(M',L')$ 的同痕类。

基本类 $[L]$ 的庞加莱对偶
$PD([L])\in H^0(L;\Lambda_{0,\mathrm{nov}})$ 是这个滤过 $A_\infty$ 代数的幺元，
且同态 $\psi_*$ 保幺。
\end{theorem}

按照定义，滤过 $A_\infty$ 代数的同构，是指具有逆的滤过 $A_\infty$ 同态。
滤过 $A_\infty$ 同态的概念定义于定义 3.2.29；它们之间的同伦定义于定义 4.2.35。
幺元和滤过 $A_\infty$ 同态保幺性的定义见定义 3.2.20 和 3.3.11。
定理 A 的证明在第 5.4 节（推论 5.4.6）完成，\emph{其中用到了第 7、8 章的结果}。

全书中相当一部分严格构造都建立在第 7、8 章的结果之上；这两章分别处理横截性和定向问题。
因此，下文不再反复写出上述斜体说明。

如果略去非平凡伪全纯圆盘的贡献，定理 A 中的滤过 $A_\infty$ 代数将与德拉姆复形同伦等价
（定理 X）。

要指出的是，定理 A 中的滤过 $A_\infty$ 结构对任意相对 spin 拉格朗日子流形 $L$ 都有定义。\liftIndexRangeStart{I008}
它在自由（形式）余代数
\[
  BH(L;\Lambda_{0,\mathrm{nov}})[1]
  =\bigoplus_k
  \underbrace{H(L;\Lambda_{0,\mathrm{nov}})[1]
  \widehat\otimes\cdots\widehat\otimes
  H(L;\Lambda_{0,\mathrm{nov}})[1]}_{k}
\]
上诱导一个余导子 $\widehat d$（见第 3.2 节公式 (3.2.1)），而 $A_\infty$ 关系
\eqref{eq:1-4} 等价于 $\widehat d\,\widehat d=0$。由此得到微分分次余代数
$(BH(L;\Lambda_{0,\mathrm{nov}})[1],\widehat d)$。它是微分分次代数经典 bar 解消的
$A_\infty$ 类比（见 [Mac63] 第 3 章第 3.5 节）。我们称它为 bar 复形。

例如，[GeJo90] 研究了 $A_\infty$ 代数的 bar 复形。与 [GeJo90] 的记号相比，我们的
$BH(L;\Lambda_{0,\mathrm{nov}})[1]$ 对应于 [GeJo90] 中的
$BH(L;\Lambda_{0,\mathrm{nov}})$。

一般而言，bar 复形的对偶是一个微分分次代数。尤里·切卡诺夫（Yuri Chekanov）曾为
$\mathbb R^3$ 中的勒让德纽结定义过一个类似的微分分次代数 [Chek02]。（在弗洛尔理论中
使用这种 bar 复形的想法，是切卡诺夫与本书作者在 20 世纪 90 年代分别独立发现的；另见
[EGH00]。）定理 A 的后半部分蕴含 bar 复形
$(BH(L;\Lambda_{0,\mathrm{nov}})[1],\widehat d)$ 的上同调是 $L$ 的不变量。然而，这个上同调
同构于 $\Lambda_{0,\mathrm{nov}}$，因为它是含幺微分分次代数的 bar 复形的 $A_\infty$ 类比
（例如见 [Mac63] 第 306 页）。

bar 复形的某些约化版本和（或）对称化版本，例如约化版本或循环约化版本，具有非平凡上同调
（关于它与环路空间的关系，见 [Fuk05II]）。定理 A 还蕴含，这些版本的上同调是相对 spin
拉格朗日子流形的不变量。

\liftIndexRangeEnd{I008}不过，我们不研究这些约化版本，而研究另一个同样由 $A_\infty$ 代数自然诱导、定义将在下文
说明的上同调。我们称它为弗洛尔上同调。这样做有几个理由。第一，bar 复形各种版本的上同调
一旦非平凡，通常是无限维的，因而较难应用；例如，其秩并不是一个好的不变量。第二，弗洛尔
上同调更直接地联系着拉格朗日子流形的各种几何性质，例如拉格朗日子流形的经典上同调，或
两个拉格朗日子流形之间的交点数。第三，在同调镜像对称对应之下，与层上同调（或 Ext 群）
对应的是弗洛尔上同调，而不是 bar 复形的上同调（见猜想 R）。

为了从定理 A 得到 $L$ 的弗洛尔上同调，作如下处理。设 $(C,\mathfrak m_k)$ 是一个滤过
$A_\infty$ 代数。注意，我们假设 $\mathfrak m_{0,0}=0$，但并不假设当 $i>0$ 时
$\mathfrak m_{0,i}\colon\mathbb Q\to C[1]$ 为零。这与一般弗洛尔理论中
$\delta\delta\neq0$ 的现象有关：具体说，\eqref{eq:1-4} 在 $k=2$ 时蕴含
\[
  \mathfrak m_1\circ\mathfrak m_1(x)
  =(-1)^{\deg x+1}\mathfrak m_2(x,\mathfrak m_0(1))
  +\mathfrak m_2(\mathfrak m_0(1),x).
\]

要得到一个上边界算子 $\delta\colon C\to C$，即满足 $\delta\delta=0$ 的映射，有如下三个层次：
\begin{enumerate}[label=(\arabic*),leftmargin=4em]
  \item $\mathfrak m_0(1)=0$ 的情形；
  \item $\mathfrak m_0(1)$ 是基本类 $PD[L]$ 的倍数的情形；
  \item 能把 $\mathfrak m$ 形变为 $\mathfrak m^{(\mathbf b_1,\mathbf b_0)}$，并使相应的
  $\mathfrak m_0^{(\mathbf b_1,\mathbf b_0)}$ 满足 (1) 或 (2) 的情形。
\end{enumerate}
在第一种情形中，显然，由 $\delta=\mathfrak m_1$ 定义的算子
$\delta\colon C\to C$ 满足 $\delta\delta=0$，因而定义一个上边界映射。在第二种情形中，可以
检验上式右端两项彼此抵消，所以仍有 $\delta=\mathfrak m_1$ 满足 $\delta\delta=0$（这一抵消论证
见补遗 [Oh93]、[ChOh03] 以及第 3.6.3 小节）。在情形 (3) 中，如果 $A_\infty$ 代数
$(C,\mathfrak m)$ 能形变到情形 (1)（分别地，形变到情形 (2)），则称它是无障碍的
（分别地，弱无障碍的）。

现在说明在无障碍情形中如何形变 $\delta$。此时考虑如下形变算子
$\delta_{\mathbf b_1,\mathbf b_0}$：
\begin{equation}
  \delta_{\mathbf b_1,\mathbf b_0}(x)
  =\sum_{k=0}^{\infty}\sum_{\ell=0}^{\infty}
  \mathfrak m_{k+\ell+1}
  \bigl(\underbrace{\mathbf b_1,\ldots,\mathbf b_1}_{k},x,
  \underbrace{\mathbf b_0,\ldots,\mathbf b_0}_{\ell}\bigr).
  \tag{1.7}\label{eq:1-7}
\end{equation}
可以证明（见第 3.8 节），若
\begin{equation}
  \sum_{k=0}^{\infty}
  \mathfrak m_k(\underbrace{\mathbf b_j,\ldots,\mathbf b_j}_{k})=0,
  \qquad j=0,1,
  \tag{1.8}\label{eq:1-8}
\end{equation}
则 $\delta_{\mathbf b_1,\mathbf b_0}\circ\delta_{\mathbf b_1,\mathbf b_0}=0$。若 $\mathbf b$ 满足~\eqref{eq:1-8}，就称
$\mathbf b$ 为定界余链\liftIndex{I004}\liftIndex{I009}。

还可以用来自环境辛流形的上同调类形变 $A_\infty$ 代数 $(C,\mathfrak m)$（见第 3.8 节）。
我们称这种形变为 bulk 形变。把上述形变论证推广到 bulk 形变的情形，便得到下面的定理 B。
这里
\[
  \Lambda^+_{0,\mathrm{nov}}
  =\left\{
    \left.\sum_{i=0}^{\infty}a_iT^{\lambda_i}e^{n_i}
    \in\Lambda_{0,\mathrm{nov}}\ \right|\ \lambda_i>0
  \right\}.
\]
参见第 1 章末尾的 (Conv.4)。

\begin{theorem}{B}\label{thm:B}
对每个相对 spin 拉格朗日子流形 $L\subset M$，可以赋予一个集合
$\mathcal M_{\mathrm{weak,def}}(L)$ 以及映射
\[
  \pi_{\mathrm{amb}}\colon\mathcal M_{\mathrm{weak,def}}(L)
  \longrightarrow H^2(M;\Lambda^+_{0,\mathrm{nov}}),
  \qquad
  \mathfrak{PO}\colon\mathcal M_{\mathrm{weak,def}}(L)
  \longrightarrow\Lambda^+_{0,\mathrm{nov}},
\]
它们具有下列性质。此后约定
\begin{equation}
\begin{aligned}
  &\mathcal M_{\mathrm{weak,def}}(L)
  \mathbin{\times_{(\pi_{\mathrm{amb}},\mathfrak{PO})}}
  \mathcal M_{\mathrm{weak,def}}(L)\\
  &\quad=\{(\mathbf b_1,\mathbf b_0)\mid
  \pi_{\mathrm{amb}}(\mathbf b_1)=\pi_{\mathrm{amb}}(\mathbf b_0),\
  \ \mathfrak{PO}(\mathbf b_1)=\mathfrak{PO}(\mathbf b_0)\}.
\end{aligned}
\tag{1.9}\label{eq:1-9}
\end{equation}

\textup{(B.1)} 对每个
\[
  (\mathbf b_1,\mathbf b_0)\in
  \mathcal M_{\mathrm{weak,def}}(L)
  \mathbin{\times_{(\pi_{\mathrm{amb}},\mathfrak{PO})}}
  \mathcal M_{\mathrm{weak,def}}(L),
\]
都有一个以 $(\mathbf b_1,\mathbf b_0)$ 为参数的弗洛尔上同调
$HF((L,\mathbf b_1),(L,\mathbf b_0);\Lambda_{0,\mathrm{nov}})$。

\textup{(B.2)} 存在乘积结构
\[
\begin{aligned}
  \mathfrak m_2\colon
  &HF((L,\mathbf b_2),(L,\mathbf b_1);\Lambda_{0,\mathrm{nov}})
  \otimes HF((L,\mathbf b_1),(L,\mathbf b_0);\Lambda_{0,\mathrm{nov}})\\
  &\longrightarrow HF((L,\mathbf b_2),(L,\mathbf b_0);\Lambda_{0,\mathrm{nov}}),
\end{aligned}
\]
其中 $(\mathbf b_1,\mathbf b_0)$、$(\mathbf b_2,\mathbf b_1)$ 均属于
$\mathcal M_{\mathrm{weak,def}}(L)
\mathbin{\times_{(\pi_{\mathrm{amb}},\mathfrak{PO})}}
\mathcal M_{\mathrm{weak,def}}(L)$。乘积 $\mathfrak m_2$ 满足结合律。特别地，对每个
$\mathbf b\in\mathcal M_{\mathrm{weak,def}}(L)$，$HF((L,\mathbf b),(L,\mathbf b);\Lambda_{0,\mathrm{nov}})$
具有环结构。

\textup{(B.3)} 若 $\psi\colon(M,L)\to(M',L')$ 是辛微分同胚，则它诱导双射
\[
  \psi_*\colon\mathcal M_{\mathrm{weak,def}}(L)
  \longrightarrow\mathcal M_{\mathrm{weak,def}}(L'),
\]
满足
\[
  \pi_{\mathrm{amb}}\circ\psi_*=\psi_*\circ\pi_{\mathrm{amb}},
  \qquad
  \mathfrak{PO}\circ\psi_*=\mathfrak{PO}.
\]
映射 $\psi_*$ 只依赖于辛微分同胚 $\psi$ 的同痕类。右端的
$\psi_*=(\psi^{-1})^*$。

\textup{(B.4)} 在 (B.3) 的情形中，有同构
\[
  \psi_*\colon HF((L,\mathbf b_1),(L,\mathbf b_0);\Lambda_{0,\mathrm{nov}})
  \xrightarrow{\ \cong\ }
  HF((L',\psi_*(\mathbf b_1)),(L',\psi_*(\mathbf b_0));\Lambda_{0,\mathrm{nov}}).
\]
而且 $\psi_*$ 与 $\mathfrak m_2$ 可交换。

\textup{(B.5)} (B.4) 中的同构 $\psi_*$ 只依赖于辛微分同胚
$\psi\colon(M,L)\to(M',L')$ 的同痕类。此外，
$(\psi\circ\psi')_*=\psi_*\circ\psi'_* $。
\end{theorem}

$\Pi_{\mathrm{amb}}$ 见 (4.3.27)，$\mathfrak{PO}$ 见引理 4.3.28。弗洛尔上同调的构造以及
定理 B 中 (B.2) 的证明在第 3.8.10 小节完成。(B.1)，特别是弗洛尔上同调只依赖于对子
$(\mathbf b_1,\mathbf b_0)$ 这一断言，来自命题 4.3.29；该命题又是第 5.2.4 小节中推论 5.2.40 的结果。
(B.3)--(B.5) 的证明在第 4.6.5 小节完成。定理 B 由定理 A（以及它的 ``def'' 版本）通过纯代数
论证推出。（不过，本书先给出定理 B 的证明，之后才完成定理 A 的证明。）

我们还考虑 $\mathcal M_{\mathrm{weak,def}}(L)$ 的如下子集：
\[
\begin{aligned}
  \mathcal M_{\mathrm{weak}}(L)
    &=\operatorname{Zero}(\mathfrak{PO}),
  &\mathcal M_{\mathrm{def}}(L)
    &=\operatorname{Zero}(\pi_{\mathrm{amb}}),\\
  \mathcal M(L)
    &=\operatorname{Zero}(\pi_{\mathrm{amb}})
      \cap\operatorname{Zero}(\mathfrak{PO}).
\end{aligned}
\]
于是，$\mathcal M(L)$ 可与定界余链 $\mathbf b$，即~\eqref{eq:1-8} 的解的规范等价类集合等同
（规范等价的定义见第 4.3 节定义 4.3.1）。正如定理 M 将要说明的，这几个集合都可以视为形式
幂级数零点集的商。对 $\mathbf b_1,\mathbf b_0\in\mathcal M(L)$，弗洛尔上同调记为
$HF((L,\mathbf b_1),(L,\mathbf b_0);\Lambda_{0,\mathrm{nov}})$；它是由~\eqref{eq:1-7} 定义的上边界算子
$\delta_{\mathbf b_1,\mathbf b_0}$ 的上同调群。对
\[
  (\mathbf b_1,\mathbf b_0)\in
  \mathcal M_{\mathrm{weak,def}}(L)
  \mathbin{\times_{(\pi_{\mathrm{amb}},\mathfrak{PO})}}
  \mathcal M_{\mathrm{weak,def}}(L)
\]
定义的 $HF((L,\mathbf b_1),(L,\mathbf b_0);\Lambda_{0,\mathrm{nov}})$ 是它的推广。更确切地说，把弗洛尔上同调
的定义从 $\mathcal M(L)$ 推广到 $\mathcal M_{\mathrm{weak}}(L)$，需要研究滤过
$A_\infty$ 代数的幺元（见第 3.3、7.3 节和第 3.7.3 小节）；从 $\mathcal M(L)$ 推广到
$\mathcal M_{\mathrm{def}}(L)$，则需要研究定理 A 中滤过 $A_\infty$ 代数的 bulk 形变\liftIndex{I010}
（见第 3.8 节）。

定理 B 蕴含，当 $\mathcal M_{\mathrm{weak,def}}(L)\neq\varnothing$ 时，弗洛尔上同调有定义。
不过应当指出，集合 $\mathcal M_{\mathrm{weak,def}}(L)$ 可能为空。它非空的一个充分条件可以用
$L$ 的上同调群来叙述。

\begin{theorem}{C}\label{thm:C}
存在一列正整数 $m_k<\dim L/2$ 以及类
\[
  [o_k^{2m_k}(L;\mathrm{weak,def})]
  \in
  \frac{H^{2m_k}(L;\mathbb Q)}
  {\operatorname{Im}\bigl(i^*\colon H^{2m_k}(M;\mathbb Q)
  \to H^{2m_k}(L;\mathbb Q)\bigr)},
  \qquad k=1,2,\ldots,
\]
使得：如果所有障碍类 $[o_k^{2m_k}(L;\mathrm{weak,def})]$ 都为零，那么
$\mathcal M_{\mathrm{weak,def}}(L)$ 非空。数 $2-2m_k$ 是有限多个同伦类的马斯洛夫指标之和；
这些同伦类属于 $\pi_2(M,L)$，并由关于 $M$ 上某个给定殆复结构的伪全纯圆盘实现。
\end{theorem}

定理 C 是定理 3.8.41、3.8.50 和 3.8.96 的推论；这些定理在第 3.8.5、3.8.9 小节中证明。
定理 C 的证明在第 3.8.10 小节完成。

\begin{remark}{1.10}\label{rem:1-10}
类 $[o_k^{2m_k}(L;\mathrm{weak,def})]$ 一般依赖于各种选择，更不必说还依赖于给定的殆复结构，
所以它们并不是 $L$ 本身的不变量。定理 M 将以更具函子性（或不变性）的方式描述集合
$\mathcal M_{\mathrm{weak}}(L)$。
\end{remark}

下面说明在一般情形中应当如何修改 (1.2.3)。此后以 $PD$ 表示庞加莱对偶。

\begin{theorem}{D}\label{thm:D}
对每个
\[
  (\mathbf b_1,\mathbf b_0)\in
  \mathcal M_{\mathrm{weak,def}}(L)
  \mathbin{\times_{(\pi_{\mathrm{amb}},\mathfrak{PO})}}
  \mathcal M_{\mathrm{weak,def}}(L),
\]
都存在一个具有下列性质的谱序列\liftIndex{I005}\liftIndex{I050}。

\textup{(D.1)}
\[
  E_2^{p,q}
  =\bigoplus_k H^k(L;\mathbb Q)\otimes
  \left(
    T^{q\lambda}\Lambda_{0,\mathrm{nov}}
    \big/ T^{(q+1)\lambda}\Lambda_{0,\mathrm{nov}}
  \right)^{(p-k)}.
\]
这里 $\lambda>0$。

\textup{(D.2)} 弗洛尔上同调
\[
  HF((L,\mathbf b_1),(L,\mathbf b_0);\Lambda_{0,\mathrm{nov}})
\]
上存在滤过
\[
  \mathfrak F^*HF((L,\mathbf b_1),(L,\mathbf b_0);\Lambda_{0,\mathrm{nov}}),
\]
使得
\[
  E_\infty^{p,q}\cong
  \frac{\mathfrak F^qHF^p((L,\mathbf b_1),(L,\mathbf b_0);\Lambda_{0,\mathrm{nov}})}
  {\mathfrak F^{q+1}HF^p((L,\mathbf b_1),(L,\mathbf b_0);\Lambda_{0,\mathrm{nov}})}.
\]

\textup{(D.3)} 考虑如下定义的子群
$K_r\subset E_r=\bigoplus_{p,q}E_r^{p,q}$：
\[
\begin{aligned}
  K_2^{p,q}
  ={}&\bigoplus_k
  PD\!\left(\ker\bigl(H_{n-k}(L;\mathbb Q)
  \longrightarrow H_{n-k}(M;\mathbb Q)\bigr)\right)\\
  &\quad\otimes
  \left(
    T^{q\lambda}\Lambda_{0,\mathrm{nov}}
    \big/T^{(q+1)\lambda}\Lambda_{0,\mathrm{nov}}
  \right)^{(p-k)},\\
  K_{r+1}
  ={}&\frac{K_r\cap\ker\delta_r}{K_r\cap\operatorname{Im}\delta_r}
  \subset E_{r+1}.
\end{aligned}
\]
再加上 $\mathbf b_0=\mathbf b_1$ 这一假设后，对每个 $r$ 都有
$\operatorname{Im}\delta_r\subseteq K_r$。

特别地，如果包含映射诱导的同态
$i_*\colon H_*(L;\mathbb Q)\to H_*(M;\mathbb Q)$ 是单射，那么谱序列在 $E_2$ 页退化\liftIndex{I029}。

\textup{(D.4)} 谱序列与 (B.2) 中的环结构相容。换言之，每个 $E_r$ 都具有一个由
$\mathfrak m_2$ 给出的环结构，并满足
\[
  \delta_r(\mathfrak m_2(x,y))
  =-\mathfrak m_2(\delta_r(x),y)
  +(-1)^{\deg x}\mathfrak m_2(x,\delta_r(y)).
\]
滤过 $\mathfrak F$ 与 (B.2) 中的环结构相容；(D.1)、(D.2) 中的同构都是环同构。
\end{theorem}

定理 D 的证明在第 6.4.5 小节完成。

在实际应用中，知道弗洛尔上同调何时不消失非常重要；这会给出某种适当形式的伪全纯曲线的
存在性定理。为此，常可利用点类 $PD[\mathrm{pt}]\in H^n(L;\mathbb Q)$ 或基本类
$PD[L]\in H^0(L;\mathbb Q)$；至少在经典上同调中，我们知道它们不为零。

\begin{theorem}{E}\label{thm:E}
在定理 D 的情形中，假设 $\mathbf b_0=\mathbf b_1$。则存在上同调类\translatornote{原书此处作 $H^n$；依紧接的同余式以及本书定理 6.1.9 的完整陈述，改为 $H^0$。}
$PD[L]'\in H^0(L;\Lambda_{0,\mathrm{nov}})$，满足
$PD[L]'\equiv PD[L]\mod\Lambda^+_{0,\mathrm{nov}}$，并具有下列性质。

\textup{(E.1)} 对每个 $r\in\mathbb Z_{\geq0}$，都有
$\delta_r(PD[L]')=0$ 且 $PD[\mathrm{pt}]\notin\operatorname{Im}(\delta_r)$；这里
$\delta_r$ 是定理 D 中谱序列的微分。

\textup{(E.2)} 如果所有以 $L$ 为边界的伪全纯圆盘的马斯洛夫指标均非正，则
$\delta_r(PD[\mathrm{pt}])=0$ 且 $PD[L]'\notin\operatorname{Im}(\delta_r)$。
\end{theorem}

定理 E 就是定理 6.1.9，其证明见第 6 章第 6.4.3 小节。定理 E 在
$\Lambda_{\mathrm{nov}}$ 系数上成立（$\Lambda_{\mathrm{nov}}$ 定义于第 1 章末尾的 (Conv.3)）。

\begin{remark}{1.11}\label{rem:1-11}
定理 E 中 $PD[L]'$ 的性质与下述事实密切相关：定理 A 中的滤过 $A_\infty$ 代数以
$PD[L]$ 为幺元（幺元的定义见第 3.2 节条件 3.2.18 和定义 3.2.20）。关于
$PD[\mathrm{pt}]$ 的断言与映射 $H_0(L;\mathbb Q)\to H_0(M;\mathbb Q)$ 是单射这一事实有关，
也与第 3.8 节所讨论的算子 $\mathfrak p$ 有关。
\end{remark}

下面讨论一对拉格朗日子流形的弗洛尔上同调\liftIndex{I015}。设 $(L_1,L_0)$ 是一对相对 spin 拉格朗日
子流形，并假设二者干净相交。也就是说，假设当 $x\in L_1\cap L_0$ 时，
$\dim(T_xL_1\cap T_xL_0)$ 是 $x$ 的局部常值函数。（这是弗洛尔理论的博特--莫尔斯理论
版本 [Bot59]。）令
\[
  \coprod_hR_h=L_1\cap L_0
\]
为其连通分支分解。在第 3.7.5 小节和第 8.8 节中，我们将由同态
\[
  \pi_1(R_h)\longrightarrow\{\pm1\}=\operatorname{Aut}(\mathbb Z)
  \subset\operatorname{Aut}(\mathbb Q)
\]
定义 $R_h$ 上的局部系统 $\Theta^-_{R_h}$。令
\[
  C(L_1,L_0;\Lambda_{0,\mathrm{nov}})
  =\bigoplus_h H(R_h;\Theta^-_{R_h})[\mu_L(R_h)]
  \otimes\Lambda_{0,\mathrm{nov}}.
\]
这里 $\mu_L(R_h)$ 是马斯洛夫指标的博特--莫尔斯版本\liftIndex{I023}，定义于命题 3.7.59 和 (3.7.60.1)。
（注意，$\mu_L([h,w])=\mu_L(R_h)$，左端定义于命题 3.7.59。）事实上，为了定义马斯洛夫指标
的博特--莫尔斯版本，还需要固定一项额外数据，在该处记为 $w$。$[\mathord\cdot]$ 表示次数移位
（即 $C[k]^d=C^{k+d}$）。此处出现次数移位的原因见第 8.8 节。

\begin{theorem}{F}\label{thm:F}
$C(L_1,L_0;\Lambda_{0,\mathrm{nov}})$ 具有以
\[
  \bigl((H(L_1;\Lambda_{0,\mathrm{nov}}),\mathfrak m_*),
  (H(L_0;\Lambda_{0,\mathrm{nov}}),\mathfrak m_*)\bigr)
\]
为代数对的含幺滤过 $A_\infty$ 双模结构。该结构在含幺滤过 $A_\infty$ 双模同构意义下良定义。
\end{theorem}

滤过 $A_\infty$ 双模的定义见第 3.7 节。$A_\infty$ 双模的同构，是指具有逆的滤过
$A_\infty$ 双模同态；滤过 $A_\infty$ 双模同态定义于第 3.7.2 小节。定理 F 的证明在第 5.4 节
（推论 5.4.21）完成。

利用定理 F，可以定义一对拉格朗日子流形的弗洛尔上同调，并得到下面的定理 G。

\begin{theorem}{G}\label{thm:G}
设 $(L_1,L_0)$ 是 $M$ 中的一对相对 spin 拉格朗日子流形，并且
\[
  (\mathbf b_1,\mathbf b_0)\in
  \mathcal M_{\mathrm{weak,def}}(L_1)
  \mathbin{\times_{(\pi_{\mathrm{amb}},\mathfrak{PO})}}
  \mathcal M_{\mathrm{weak,def}}(L_0)
\]
（右端的定义见~\eqref{eq:1-9}）。那么可以赋予一个具有 (G.1)--(G.5) 所列性质的弗洛尔
上同调 $HF((L_1,\mathbf b_1),(L_0,\mathbf b_0);\Lambda_{0,\mathrm{nov}})$。令
\[
  HF((L_1,\mathbf b_1),(L_0,\mathbf b_0);\Lambda_{0,\mathrm{nov}})
  \otimes_{\Lambda_{0,\mathrm{nov}}}\Lambda_{\mathrm{nov}}
  =HF((L_1,\mathbf b_1),(L_0,\mathbf b_0);\Lambda_{\mathrm{nov}}).
\]

\textup{(G.1)} 如果 $L_0=L_1=L$，那么
$HF((L_1,\mathbf b_1),(L_0,\mathbf b_0);\Lambda_{0,\mathrm{nov}})$ 与定理 B 中的弗洛尔上同调一致。

\textup{(G.2)} 如果 $L_1$ 与 $L_0$ 干净相交，则
\[
  \operatorname{rank}_{\Lambda_{\mathrm{nov}}}
  HF((L_1,\mathbf b_1),(L_0,\mathbf b_0);\Lambda_{\mathrm{nov}})
  \leq\sum_{h,k}\operatorname{rank}_{\mathbb Q}
  H^k(R_h;\Theta^-_{R_h}),
\]
其中每个 $R_h$ 都是 $L_0\cap L_1$ 的一个连通分支，$\Theta^-_{R_h}$ 是其上的局部系统，
它与某个表示
$\pi_1(R_h)\to\operatorname{Aut}(\mathbb R)=\{\pm1\}$ 相伴。

\textup{(G.3)} 若 $\psi\colon M\to M'$ 是满足
$\psi(L_i)=L_i'$（$i=0,1$）的辛微分同胚，则有典范同构
\[
  \psi_*\colon
  HF((L_1,\mathbf b_1),(L_0,\mathbf b_0);\Lambda_{0,\mathrm{nov}})
  \xrightarrow{\ \cong\ }
  HF((L_1',\psi_*\mathbf b_1),(L_0',\psi_*\mathbf b_0);\Lambda_{0,\mathrm{nov}}),
\]
其中 $\psi_*\colon\mathcal M_{\mathrm{weak,def}}(L_i)\to
\mathcal M_{\mathrm{weak,def}}(L_i')$ 与定理 B 中相同。同构 $\psi_*$ 只依赖于满足
$\psi(L_i)=L_i'$ 的辛微分同胚 $\psi$ 的同痕类，并且
$(\psi\circ\psi')_*=\psi_*\circ\psi'_* $。

\textup{(G.4)} 若
$\psi_i^s\colon M\to M$（$i=0,1$，$s\in[0,1]$）是哈密顿同痕，满足
$\psi_i^0=\operatorname{identity}$、$\psi_i^1(L_i)=L_i'$，则它诱导同构
\[
\begin{aligned}
  (\psi_1^s,\psi_0^s)_*\colon
  &HF((L_1,\mathbf b_1),(L_0,\mathbf b_0);\Lambda_{\mathrm{nov}})\\
  &\xrightarrow{\ \cong\ }
  HF((L_1',\psi_{1*}^1\mathbf b_1),(L_0',\psi_{0*}^1\mathbf b_0);\Lambda_{\mathrm{nov}}).
\end{aligned}
\]
同构 $(\psi_1^s,\psi_0^s)_*$ 只依赖于满足上述条件的哈密顿同痕
$\psi_i^s$（$i=0,1$，$s\in[0,1]$）的同痕类；对于哈密顿同痕的复合，该同构具有函子性。

\textup{(G.5)} $HF((L_1,\mathbf b_1),(L_0,\mathbf b_0);\Lambda_{0,\mathrm{nov}})$ 是下列环对上的双模：
\[
  \bigl(
    HF((L_1,\mathbf b_1),(L_1,\mathbf b_1);\Lambda_{0,\mathrm{nov}}),
    HF((L_0,\mathbf b_0),(L_0,\mathbf b_0);\Lambda_{0,\mathrm{nov}})
  \bigr).
\]
(G.3)、(G.4) 中的同构都是双模同构。
\end{theorem}

\begin{remark}{1.12}\label{rem:1-12}
上述 (G.3) 断言，当我们用同一个辛微分同胚移动 $L_1,L_0$ 时，弗洛尔上同调仍然同构。
(G.4) 断言，即使用不同的哈密顿同痕彼此独立地移动 $L_1,L_0$，弗洛尔上同调也不改变。
后一种不变性只在 $\Lambda_{\mathrm{nov}}$ 系数上成立，在 $\Lambda_{0,\mathrm{nov}}$ 系数上并不成立。
定理 J 将进一步讨论这一点。
\end{remark}

弗洛尔上同调的构造以及定理 G 中 (G.1)、(G.5) 的证明在第 3.8.10 小节完成。(G.3) 是
定理 4.1.4、5.3.64 的推论，这两个定理分别在第 5.3.1、5.3.6 小节中证明。(G.4) 是定理
4.1.5、5.3.65 的推论，这两个定理分别在第 5.3.2--5.3.6 小节中证明。(G.3)、(G.4) 的证明
在第 5.3.7 小节完成。(G.2) 在第 5.4.6 小节中证明，也可以由第 6 章中证明的定理 6.1.4 推出。
对非干净相交的情形，证明定理 G 还需要更多论证；这些论证见第 6.5.4 小节（命题 6.5.38、
定义 6.5.39）。

定理 G 也可由定理 F 所给滤过 $A_\infty$ 双模的函子性推出（见定理 5.3.1、5.3.14）。
（不过，本书在完成定理 F 的证明之前先给出定理 G 的证明。）

还要指出，定理 F 蕴含：对任意相对 spin 对 $(L_1,L_0)$，都可以构造各种 bar 复形，且其
上同调与各种选择无关。构造这样的 bar 复形不需要假设
$\mathcal M(L_1)\neq\varnothing$、$\mathcal M(L_0)\neq\varnothing$（见定义 3.7.5）。
我们用 $\mathcal M(L_j)$ 中的元素把 bar 复形缩减到较为合理并且适于应用的规模。

对于干净相交的一对拉格朗日子流形 $(L_1,L_0)$，也有一个谱序列（见定理 6.1.4）。

还可以把这一点推广到涉及三个或更多拉格朗日子流形的乘积结构。具体说，我们得到一个滤过
$A_\infty$ 范畴。这种范畴的存在性是第一位署名作者深谷贤治（Kenji Fukaya）约在 1993 年
注意到的，其灵感来自
唐纳森 1992 年在华威大学所作的演讲。其构造细节见
[Fuk02II]；该文以本书内容为前提。
